\documentclass[11pt,a4paper]{article}
\usepackage[margin=1in]{geometry}
\usepackage{amsmath,amssymb,amsthm,bm}
\usepackage{mathtools}
\usepackage{hyperref}
\usepackage{booktabs}

\title{Three-layer erf networks: equations of state in Hermite matrix elements\\
\large with a non-Gaussian hidden field}
\author{Akiva Lipshitz}
\date{\today}

\newtheorem{remark}{Remark}
\newcommand{\E}{\mathbb{E}}
\newcommand{\erf}{\operatorname{erf}}
\newcommand{\tr}{\operatorname{tr}}
\newcommand{\He}{\mathrm{He}}

\begin{document}
\maketitle

\begin{abstract}
We derive mean-field equations of state for the Bayesian posterior of a three-layer erf network trained on a target
$\He_1+\varepsilon\He_3$.
We work in the subspace spanned by the odd Hermite polynomials $\He_1,\He_3,\dots,\He_{n_{\max}}$ of the target direction, and keep every kernel as a full matrix in this subspace.
The second hidden layer's field $h$ keeps its exact, non-Gaussian distribution, and the read-in weights keep their exact distribution
in the two coordinates (aligned component, perpendicular norm) that the equations depend on.
The term that the hidden layer adds to the action of each read-in row has the form $J^{-1}(J-H)J^{-1}$ whether or not $h$ is Gaussian.
Making every kernel diagonal and linearising the hidden nonlinearity reduces the system to three scalar equations.
Finally, we show why the equations at infinite first-layer width do not reduce to those of a two-layer network.
\end{abstract}

%% ================================================================
\section{Model and notation}
%% ================================================================
\paragraph{Network.}
\begin{equation}
  f(x)=\sum_{k=1}^{N_2}a_k\,\erf\big(h_k(x)\big),\qquad
  h_k(x)=\sum_{j=1}^{N_1}B_{kj}\,\erf(w_j\cdot x),\qquad x\in\mathbb R^d .
  \label{eq:net}
\end{equation}
The inputs are $x\sim\mathcal N(0,I_d)$.
The target is $y(x)=\He_1(x_0)+\varepsilon\He_3(x_0)$, where $x_0=v_*\cdot x$ for a unit vector $v_*$, and
$\He_3(z)=(z^3-3z)/\sqrt6$.
We write $\Phi=\erf$ when it acts as an activation.

\paragraph{Posterior.}
We use the dataset-averaged form of the likelihood, in which the sum over $P$ training points is replaced by $P$ times the population
average. The posterior is then $p\propto e^{-S_0}$ with
\begin{equation}
  S_0=\frac{N_2\chi}{2\sigma_A^2}\|a\|^2+\frac{N_1}{2\sigma_B^2}\|B\|^2+\frac{d}{2\sigma_w^2}\|W\|^2
  +\frac{P\chi}{2\kappa}\,\|f-y\|_x^2 .
  \label{eq:S0}
\end{equation}

\paragraph{Inner products.}
We write $\langle g\mid h\rangle_x=\int d\mu_x\,g(x)h(x)$ and $\|g\|_x^2=\langle g|g\rangle_x$, with $d\mu_x$ the Gaussian input measure.
A kernel acts as $\langle g\mid A\mid h\rangle=\iint d\mu_xd\mu_{x'}\,g(x)A(x,x')h(x')$.

\paragraph{The Hermite subspace.}
Let $\varphi_n(x)=\He_n(x_0)/\sqrt{n!}$.
These are orthonormal, $\langle\varphi_m|\varphi_n\rangle_x=\delta_{mn}$.
We work in the subspace of functions
\begin{equation*}
  \mathrm{span}\{\varphi_1,\varphi_3,\dots,\varphi_{n_{\max}}\}
  =\mathrm{span}\{\He_1(x_0),\He_3(x_0),\dots,\He_{n_{\max}}(x_0)\}.
\end{equation*}
Its index set is $S=\{1,3,\dots,n_{\max}\}$ and its dimension is $s=(n_{\max}+1)/2$.
The target lies in this subspace.
The part of a function orthogonal to the subspace is marked with $\perp$.
For a kernel $A(x,x')$, the bold symbol $\mathbf A$ denotes the $s\times s$ matrix of its elements in the subspace,
\begin{equation*}
  A_{mn}=\langle\varphi_m|A|\varphi_n\rangle,\qquad m,n\in S.
\end{equation*}
All entries are kept, \emph{including} those with $m\neq n$.
Bold lower-case letters are column vectors of length $s$, with one entry per $n\in S$; for the target,
$\mathbf y=(y_1,y_3,\dots,y_{n_{\max}})^\top=(1,\varepsilon,0,\dots,0)^\top$.

\paragraph{Shorthands.}
\begin{equation}
  \beta=\frac{d}{2\sigma_w^2},\qquad \tilde g=\frac{\sigma_A^2\chi}{N_2},\qquad r_B=\sigma_B^2\,\frac{N_2}{N_1}.
\end{equation}
$\beta$ is the coefficient of $|w|^2$ in the prior of a read-in row.
$\tilde g$ sets how strongly the data term acts on each hidden field $h_k$ (Section~2, eq.~\eqref{eq:Sh}).
$r_B$ sets how strongly the $N_2$ hidden fields act back on each read-in row (eq.~\eqref{eq:oneW}):
each of the $N_2$ hidden fields couples to a given read-in row through one weight $B_{kj}$ of variance $\sigma_B^2/N_1$,
so the total is $N_2\cdot\sigma_B^2/N_1$.
It vanishes when $N_1\to\infty$ at fixed $N_2$.

%% ================================================================
\section{Field theory and mean-field decoupling}
%% ================================================================
\paragraph{Fields.}
Insert functional delta functions to define the output field $f$ and the hidden fields $h_k$, and represent them with conjugate fields
$\tau$ and $\nu_k$.
The weights $a$ and $B$ then appear linearly in the exponent, and the Gaussian integrals over them give
\begin{multline}
  S=\frac{d}{2\sigma_w^2}\|W\|^2-i\langle\tau|f\rangle-i\sum_{k}\langle\nu_k|h_k\rangle+\frac{P\chi}{2\kappa}\|f-y\|_x^2\\
  +\frac{\sigma_A^2}{2N_2\chi}\sum_{k=1}^{N_2}\langle\tau\mid\Phi\circ h_k\rangle^2
  +\frac{\sigma_B^2}{2N_1}\sum_{k=1}^{N_2}\sum_{j=1}^{N_1}\langle\nu_k\mid\Phi\circ W_j\rangle^2 ,
  \label{eq:S}
\end{multline}
where $W_j$ is row $j$ of $W$ and $(\Phi\circ W_j)(x)=\erf(w_j\cdot x)$.
The last two terms are quartic and couple the layers.

\paragraph{Mean-field decoupling.}
Each quartic term has the form $\frac12\iint X(x)X(x')\,Y(x)Y(x')$.
We replace it by $\frac12\iint\big(\E[XX]\,YY+XX\,\E[YY]\big)$.
This drops the product of the two fluctuations and a constant.
Five kernels appear:
\begin{equation}
\begin{gathered}
  J(x,x')=\E\big[\erf(w\cdot x)\erf(w\cdot x')\big],\qquad
  K(x,x')=\E\big[\erf(h(x))\erf(h(x'))\big],\\
  H(x,x')=\E[h(x)h(x')],\qquad
  \mathcal V(x,x')=\E[\nu(x)\nu(x')],\qquad
  \mathcal T(x,x')=\E[\tau(x)\tau(x')].
\end{gathered}
\label{eq:kernels}
\end{equation}
All neurons in a layer have the same distribution (the action is unchanged if they are permuted), so we drop the neuron indices.
Each expectation is taken with the decoupled action, and that action itself contains the kernels; the kernels are therefore found
by solving equations in which they appear on both sides (Section~5).
After decoupling, the action is a sum of separate terms, each containing a single read-in row, a single hidden field, or the output:
\begin{align}
  &\text{read-in $W_j$:}&& \frac{d}{2\sigma_w^2}|W_j|^2+\frac{r_B}{2}\,\langle\Phi\circ W_j\mid\mathcal V\mid\Phi\circ W_j\rangle,
  \label{eq:oneW}\\
  &\text{hidden $h_k$:}&& -i\langle\nu_k|h_k\rangle+\frac{\sigma_B^2}{2}\langle\nu_k|J|\nu_k\rangle
     +\frac{\sigma_A^2}{2N_2\chi}\langle\Phi\circ h_k\mid\mathcal T\mid\Phi\circ h_k\rangle,
  \label{eq:oneh}\\
  &\text{output:}&& -i\langle\tau|f\rangle+\frac{\sigma_A^2}{2\chi}\langle\tau|K|\tau\rangle+\frac{P\chi}{2\kappa}\|f-y\|_x^2 .
  \label{eq:onef}
\end{align}
The factor $r_B=\sigma_B^2N_2/N_1$ in \eqref{eq:oneW} comes from $\sum_{k=1}^{N_2}\nu_k\nu_k\to N_2\mathcal V$.
\emph{This is the only place where $N_1$ appears.}

\paragraph{Output.}
Integrating $f$ out of \eqref{eq:onef} leaves a Gaussian action for $\tau$,
\begin{equation*}
  -i\langle\tau|y\rangle+\tfrac1{2\chi}\langle\tau|\sigma_A^2K+\tfrac\kappa P|\tau\rangle .
\end{equation*}
Hence
\begin{equation}
  \E\tau=i\chi\,t,\qquad \big(\sigma_A^2K+\tfrac\kappa P\big)\,t=y,\qquad
  \bar f=\E f=y-\tfrac\kappa P\,t=\sigma_A^2K\,t,
  \label{eq:readout}
\end{equation}
and $\mathcal T=\chi(\sigma_A^2K+\tfrac\kappa P)^{-1}-\chi^2\,t(x)t(x')$.
We keep the leading $O(\chi^2)$ part, $\mathcal T(x,x')\simeq-\chi^2\,t(x)t(x')$.
The learnability of mode $n$ is $\mu_n=\bar f_n/y_n$.

\paragraph{Hidden field.}
In \eqref{eq:oneh}, $\nu_k$ appears only quadratically, so it can be integrated out:
\begin{equation}
  \int\!\mathcal D\nu\;\exp\!\big(i\langle\nu|h\rangle-\tfrac12\langle\nu|G|\nu\rangle\big)\;\propto\;\exp\!\big(-\tfrac12\langle h|G^{-1}|h\rangle\big),
  \qquad G=\sigma_B^2J .
  \label{eq:G}
\end{equation}
Inserting $\mathcal T(x,x')\simeq-\chi^2t(x)t(x')$, the action of one hidden field is
\begin{equation}
  S[h]=\tfrac12\langle h|G^{-1}|h\rangle-\frac{\tilde g}{2}\,\langle t\mid\erf\circ h\rangle^2 .
  \label{eq:Sh}
\end{equation}
This action is \emph{not} quadratic in $h$, and we do not approximate it as quadratic.

\paragraph{The kernel $\mathcal V$ acting on the read-in rows.}
Given $h$, the action for $\nu$ is $-i\langle\nu|h\rangle+\frac12\langle\nu|G|\nu\rangle$.
So $\nu$ is Gaussian with mean $iG^{-1}h$ and covariance $G^{-1}$, and
$\E[\nu\nu\mid h]=G^{-1}-G^{-1}h\,h^\top G^{-1}$.
The second term of \eqref{eq:Sh}, $-\frac{\tilde g}{2}\langle t|\erf\circ h\rangle^2$, depends on $h$ and not on $\nu$.
So averaging over the distribution of $h$ gives, exactly,
\begin{equation}
  \mathcal V=G^{-1}(G-H)G^{-1},\qquad H=\E[hh].
  \label{eq:V}
\end{equation}
This holds for any distribution of $h$, Gaussian or not.

%% ================================================================
\section{Reduction to Hermite matrix elements}
%% ================================================================
\paragraph{Read-in features.}
Write a read-in row as $w=(w_0,w_\perp)$, where $w_0=w\cdot v_*$ and $q=|w_\perp|^2$, and let $D_w=1+2(w_0^2+q)$.
At fixed $x_0$, $w\cdot x=w_0x_0+w_\perp\cdot x_\perp$ is exactly Gaussian over $x_\perp$, with variance $q$.
Hence $\E_{x_\perp}\erf(w\cdot x)=\erf\big(w_0x_0/\sqrt{1+2q}\big)$.
Gaussian integration by parts, $\E[g(x_0)\He_n(x_0)]=\E[g^{(n)}(x_0)]$, then gives
\begin{gather}
  c_n(w_0,q):=\langle\varphi_n\mid\erf(w\cdot x)\rangle_x
  =\frac{2}{\sqrt\pi}\,\frac{(-2)^{\frac{n-1}{2}}(n-2)!!}{\sqrt{n!}}\;\frac{w_0^{\,n}}{D_w^{n/2}},
  \label{eq:cn}\\
  c_1=\frac{2}{\sqrt\pi}\frac{w_0}{D_w^{1/2}},\qquad
  c_3=-\frac{4}{\sqrt{6\pi}}\frac{w_0^3}{D_w^{3/2}} . \nonumber
\end{gather}
The sign of $c_3$ is negative.
It cancels in $\E c_3^2$, but it fixes the sign of $\E[c_1c_3]$.
The arcsine identity $\E[\erf(u)\erf(u')]=\frac2\pi\arcsin\frac{2\,\mathrm{cov}(u,u')}{\sqrt{(1+2\,\mathrm{var}\,u)(1+2\,\mathrm{var}\,u')}}$ for jointly
Gaussian $(u,u')$, applied over $x$, gives
\begin{gather}
  \|\erf(w\cdot x)\|_x^2=\frac2\pi\arcsin\frac{2|w|^2}{1+2|w|^2}, \nonumber\\
  k(w,w'):=\langle\erf(w\cdot x)\mid\erf(w'\cdot x)\rangle_x
  =\frac2\pi\arcsin\frac{2\,w\cdot w'}{\sqrt{(1+2|w|^2)(1+2|w'|^2)}} .
  \label{eq:kww}
\end{gather}
The part of a feature outside the subspace has squared norm $P_\perp(w)=\|\erf(w\cdot x)\|_x^2-\sum_{n\in S}c_n^2$.

\paragraph{First-layer kernel.}
Collect the read-in features into the vector
\begin{equation*}
  \mathbf c(w)=\big(c_1(w),\,c_3(w),\,\dots,\,c_{n_{\max}}(w)\big)^\top\in\mathbb R^{s}.
\end{equation*}
Projecting $J$ in \eqref{eq:kernels} onto the subspace gives $J_{mn}=\E_w[c_m(w)c_n(w)]$, that is,
\begin{equation}
  \mathbf J=\E_w\big[\mathbf c\,\mathbf c^\top\big],\qquad \mathbf G=\sigma_B^2\mathbf J ,
  \label{eq:JS}
\end{equation}
where $\E_w$ is over the distribution of read-in weights, which is non-Gaussian in the feature-learning regime \eqref{eq:readin}.

\paragraph{The hidden field in the subspace.}
Expand $h(x)=\sum_{n\in S}h_n\varphi_n(x)+h_\perp(x)$ with $h_n=\langle\varphi_n|h\rangle_x$, and write
$\mathbf h=(h_1,h_3,\dots,h_{n_{\max}})^\top$.
Under the prior \eqref{eq:G}, $\mathbf h\sim\mathcal N(0,\mathbf G)$.
Here $\mathbf h$ holds the $s$ components of $h$ in the subspace and $h_\perp$ is the rest of $h$.
We take the covariance of $h_\perp$ to be $G$ with its block in the subspace removed:
\begin{equation}
  G_\perp(x,x')=G(x,x')-\sum_{m,n\in S}\varphi_m(x)\,G_{mn}\,\varphi_n(x').
  \label{eq:Gperp}
\end{equation}
\paragraph{Hermite coefficients of the hidden activation.}
Since $y$ lies in the subspace, so does $t$ in \eqref{eq:readout}, which is then the vector $\mathbf t=(t_1,t_3,\dots,t_{n_{\max}})^\top$.
The second term of \eqref{eq:Sh} then involves the activation $\erf\circ h$ only through its Hermite coefficients
\begin{equation*}
  e_n:=\langle\varphi_n\mid\erf\circ h\rangle_x,\qquad n\in S,\qquad
  \mathbf e=(e_1,e_3,\dots,e_{n_{\max}})^\top,
\end{equation*}
since $\langle t\mid\erf\circ h\rangle=\mathbf t\cdot\mathbf e$.
They play for the hidden field the role that $c_n$ plays for a read-in row.
To compute them, split $h(x)=a(x_0)+h_\perp(x)$ with $a(x_0)=\sum_{m\in S}h_m\varphi_m(x_0)$.
The remaining part $h_\perp$ enters only through one number, its squared norm
\begin{equation}
  q_h:=\|h_\perp\|_x^2=\int d\mu_x\,h_\perp(x)^2 ,
\end{equation}
because at fixed $x_0$ we treat $h_\perp(x)$ as a Gaussian variable over $x_\perp$ with variance $q_h$.
The identity $\E_\xi\erf(a+\sqrt q\,\xi)=\erf\big(a/\sqrt{1+2q}\big)$ for $\xi\sim\mathcal N(0,1)$, already used for $c_n$, then gives
\begin{equation}
  e_n(\mathbf h,q_h)
  =\E_{x_0}\!\Big[\varphi_n(x_0)\,\erf\Big(\frac{\sum_{m\in S}h_m\varphi_m(x_0)}{\sqrt{1+2q_h}}\Big)\Big],
  \qquad x_0\sim\mathcal N(0,1).
  \label{eq:Psi}
\end{equation}
This keeps the full dependence on $\mathbf h$; it is not linearised.
For $\mathbf h=(h_1,0,0,\dots)^\top$ it reduces to $e_n=c_n(h_1,q_h)$.

\paragraph{The distribution of $q_h$.}
Let $G_\perp(x,x')=\sum_\ell\lambda^\perp_\ell\,\psi_\ell(x)\psi_\ell(x')$ be the eigen-expansion of $G_\perp$.
Its eigenfunctions $\psi_\ell$ are orthonormal and orthogonal to every $\varphi_n$.
Under the prior, $h_\perp=\sum_\ell\sqrt{\lambda^\perp_\ell}\,g_\ell\psi_\ell$ with independent $g_\ell\sim\mathcal N(0,1)$, so
\begin{equation*}
  q_h=\|h_\perp\|_x^2=\sum_\ell\lambda^\perp_\ell\,g_\ell^2,\qquad
  \E q_h=\sum_\ell\lambda^\perp_\ell=\tr G_\perp,\qquad
  \operatorname{Var}q_h=2\sum_\ell(\lambda^\perp_\ell)^2=2\tr G_\perp^2 ,
\end{equation*}
using $\E g^2=1$ and $\operatorname{Var}g^2=2$.

\emph{Equal eigenvalues.}
Suppose $G_\perp$ had exactly $M$ nonzero eigenvalues, all equal to $\lambda$.
Then $q_h=\lambda\sum_{\ell=1}^Mg_\ell^2$ is $\lambda$ times a chi-squared variable with $M$ degrees of freedom, and its density is
\begin{equation*}
  p(q_h)=\frac{q_h^{M/2-1}\,e^{-q_h/(2\lambda)}}{\Gamma(M/2)\,(2\lambda)^{M/2}},
  \qquad \E q_h=M\lambda,\quad \operatorname{Var}q_h=2M\lambda^2,
\end{equation*}
where $\Gamma(\cdot)$ is the Gamma function.

\emph{General eigenvalues.}
For unequal $\lambda^\perp_\ell$, the density of $\sum_\ell\lambda^\perp_\ell g_\ell^2$ has no closed form.
We approximate it by the density of the same form, written in terms of its mean $\bar q_h=M\lambda$:
\begin{equation}
  p(q_h)\;\propto\;q_h^{M/2-1}\exp\!\Big(-\frac{M\,q_h}{2\bar q_h}\Big),
  \qquad \E q_h=\bar q_h,\quad \operatorname{Var}q_h=\frac{2\bar q_h^{\,2}}{M}.
  \label{eq:gamma}
\end{equation}
Here $M$ is allowed to be any real number $M>0$.
We fix $\bar q_h$ and $M$ by requiring the mean and the variance to match the exact ones:
\begin{equation}
  \bar q_h=\tr G_\perp,\qquad
  \frac{2\bar q_h^{\,2}}{M}=2\tr G_\perp^2
  \quad\Longrightarrow\quad
  M=\frac{(\tr G_\perp)^2}{\tr G_\perp^2}=\frac{\big(\sum_\ell\lambda^\perp_\ell\big)^2}{\sum_\ell(\lambda^\perp_\ell)^2}.
  \label{eq:M}
\end{equation}
$M$ is the effective number of directions in which $h_\perp$ fluctuates.
It equals the number of nonzero eigenvalues when they are all equal, and is smaller when they are spread out.

\paragraph{The traces.}
Using $\langle\varphi_m|\varphi_n\rangle=\delta_{mn}$ and exchanging the $x$- and $w$-averages,
\begin{align}
  \bar q_h&=\sigma_B^2\Big(\E_w\|\erf(w\cdot x)\|_x^2-\tr \mathbf J\Big),
  \label{eq:qbar}\\
  \tr G_\perp^2&=\sigma_B^4\Big(\E_{w,w'}\big[k(w,w')^2\big]-\|\mathbf J\|_F^2\Big).
  \label{eq:trG2}
\end{align}
The second line uses $\iint J(x,x')^2=\E_{w,w'}\langle\erf(w\cdot x)|\erf(w'\cdot x)\rangle_x^2$, with $w,w'$ independent draws from the distribution of read-in weights.
Writing $w\cdot w'=w_0w_0'+\sqrt{qq'}\,u$, the cosine $u$ between two independent directions in $\mathbb R^{d-1}$ has density
$\propto(1-u^2)^{(d-4)/2}$.

\begin{remark}
Equation \eqref{eq:Gperp} treats $\mathbf h$ and $h_\perp$ as independent.
The exact covariance of $h_\perp=h-\sum_{n\in S}h_n\varphi_n$ also subtracts the cross terms:
\begin{equation*}
  G(x,x')-\sum_{n\in S}\Big[\varphi_n(x)\,(G\varphi_n)(x')+(G\varphi_n)(x)\,\varphi_n(x')\Big]+\sum_{m,n\in S}\varphi_m(x)\,G_{mn}\,\varphi_n(x'),
\end{equation*}
where $(G\varphi_n)(x)=\int d\mu_{x'}\,G(x,x')\varphi_n(x')$.
It has the same trace, but its squared trace norm is smaller by
$2\sum_{n\in S}\big(\|G\varphi_n\|_x^2-\sum_{m\in S}G_{mn}^2\big)\ge0$.
So \eqref{eq:Gperp} slightly underestimates $M$; $\bar q_h$ is unaffected.
\end{remark}

%% ================================================================
\section{The distribution of each layer}
%% ================================================================
\paragraph{Hidden field.}
The action \eqref{eq:Sh} governs the hidden field $h$.
After the steps of Section~3, $h$ is described by $s+1$ numbers: its $s$ components $\mathbf h=(h_1,h_3,\dots,h_{n_{\max}})^\top$
in the subspace, and the squared norm $q_h$ of its part outside the subspace.
Combining the Gaussian prior of $\mathbf h$, the distribution \eqref{eq:gamma} of $q_h$, and the second term of \eqref{eq:Sh},
their joint distribution is
\begin{equation}
  -\log p(\mathbf h,q_h)=\tfrac12\,\mathbf h^\top \mathbf G^{-1}\mathbf h
  -\big(\tfrac M2-1\big)\log q_h+\tfrac{M}{2\bar q_h}\,q_h
  -\frac{\tilde g}{2}\,\big(\mathbf t\cdot\mathbf e(\mathbf h,q_h)\big)^2 .
  \label{eq:middle}
\end{equation}
$\mathbf e$ is bounded and nonlinear in $\mathbf h$, so the distribution is not Gaussian, and $-\log p$ can have local minima away from the origin.
All expectations in this note are over the full distributions, including any such minima.

Its moments give the matrices in the subspace of the kernels \eqref{eq:kernels}:
\begin{equation}
  \mathbf K=\E\big[\mathbf e\,\mathbf e^\top\big],\qquad
  \mathbf H=\E\big[\mathbf h\,\mathbf h^\top\big],\qquad
  \langle q_h\rangle=\E[q_h].
  \label{eq:moments}
\end{equation}

\paragraph{Read-in layer.}
Split $\Phi\circ w=\sum_{n\in S}c_n\varphi_n+(\Phi\circ w)_\perp$, and treat $\mathcal V$ as having no elements that connect the subspace to its
orthogonal complement.
In the subspace, \eqref{eq:V} gives $\boldsymbol{\mathcal V}=\mathbf G^{-1}(\mathbf G-\mathbf H)\mathbf G^{-1}$.
Outside the subspace, we again treat $G_\perp$ as having $M$ equal eigenvalues, as in \eqref{eq:gamma}: take $G_\perp\simeq\frac{\bar q_h}{M}\Pi_M$ and
$H_\perp\simeq\frac{\langle q_h\rangle}{M}\Pi_M$, where $\Pi_M$ projects onto $M$ of the $\psi_\ell$.
These have the trace $\bar q_h$ and the effective number of directions $M$ of \eqref{eq:M}.
This gives $\mathcal V_\perp\simeq\frac{M(\bar q_h-\langle q_h\rangle)}{\bar q_h^{\,2}}\Pi_M$.
In the coordinates $(w_0,q)$, with the Jacobian $q^{(d-3)/2}$ of $w_\perp\mapsto q$, the action \eqref{eq:oneW} of a read-in row becomes
\begin{multline}
  -\log p(w_0,q)=\beta(w_0^2+q)-\tfrac{d-3}{2}\log q
  +\frac{r_B}{2}\,\mathbf c^\top \mathbf G^{-1}\big(\mathbf G-\mathbf H\big)\mathbf G^{-1}\mathbf c\\
  +\frac{r_B}{2}\,\frac{M\big(\bar q_h-\langle q_h\rangle\big)}{\bar q_h^{\,2}}\,P_\perp(w)
  +\text{const}.
  \label{eq:readin}
\end{multline}
The first $r_B$ term contains $c_1^2\propto w_0^2/D_w$, $c_3^2\propto w_0^6/D_w^3$, and the cross term $c_1c_3\propto w_0^4/D_w^2$.
None of them is dropped.

\begin{remark}[Variational form of the readout]
Let $Z_h(\mathbf t)$ be the normaliser of \eqref{eq:middle}.
Then $\nabla_{\mathbf t}\log Z_h=\tilde g\,\mathbf K\mathbf t$, so the readout condition \eqref{eq:readout} is the stationarity condition of the convex potential
\begin{equation}
  \mathcal F(\mathbf t)=\frac{\kappa}{2P}|\mathbf t|^2+\frac{N_2}{\chi}\log Z_h(\mathbf t)-\mathbf y\cdot \mathbf t,
  \qquad
  \nabla^2\mathcal F=\tfrac\kappa P+\sigma_A^2\big(\mathbf K+\tilde g\,\operatorname{Cov}[(\mathbf t\cdot\mathbf e)\mathbf e]\big).
\end{equation}
\end{remark}

%% ================================================================
\section{The closed system}
%% ================================================================
The unknowns are $(\mathbf J,\bar q_h,M)$ for the read-in layer and $\mathbf t=(t_1,t_3,\dots,t_{n_{\max}})^\top$ for the readout.
The constants are $\beta=d/(2\sigma_w^2)$, $\tilde g=\sigma_A^2\chi/N_2$ and $r_B=\sigma_B^2N_2/N_1$.
The variables and functions that appear are the following.
\begin{itemize}
  \item A read-in row is $w=(w_0,w_\perp)$, with $w_0=w\cdot v_*$ its component along the target direction and
    $q=|w_\perp|^2$ the squared norm of the rest. Let $D_w=1+2(w_0^2+q)$.
  \item $\mathbf c(w_0,q)=(c_1,c_3,\dots,c_{n_{\max}})^\top$ are the Hermite coefficients of the first-layer activation,
    $c_n=\langle\varphi_n|\erf(w\cdot x)\rangle_x
    =\frac{2}{\sqrt\pi}\frac{(-2)^{(n-1)/2}(n-2)!!}{\sqrt{n!}}\frac{w_0^n}{D_w^{n/2}}$, eq.~\eqref{eq:cn}.
  \item $\mathbf h=(h_1,h_3,\dots,h_{n_{\max}})^\top$, with $h_n=\langle\varphi_n|h\rangle_x$, are the components of the hidden field
    $h$ in the subspace, and $q_h=\|h_\perp\|_x^2$ is the squared norm of its part outside the subspace.
  \item $\mathbf e(\mathbf h,q_h)=(e_1,e_3,\dots,e_{n_{\max}})^\top$ are the Hermite coefficients of the hidden activation,
    $e_n=\langle\varphi_n|\erf\circ h\rangle_x
    =\E_{x_0}\big[\varphi_n(x_0)\,\erf\big(\sum_{m\in S}h_m\varphi_m(x_0)/\sqrt{1+2q_h}\big)\big]$, eq.~\eqref{eq:Psi}.
  \item $k(w,w')=\frac2\pi\arcsin\frac{2\,w\cdot w'}{\sqrt{(1+2|w|^2)(1+2|w'|^2)}}$ is the overlap of two first-layer
    activations, eq.~\eqref{eq:kww}, and $\|\erf(w\cdot x)\|_x^2=k(w,w)$.
  \item $P_\perp(w)=k(w,w)-|\mathbf c(w)|^2$ is the squared norm of the first-layer activation outside the subspace.
\end{itemize}
\begin{align}
  &\text{(E1)}\quad \mathbf G=\sigma_B^2\mathbf J, \nonumber\\
  &\text{(E2)}\quad p(\mathbf h,q_h)\propto e^{-\frac12\mathbf h^\top\mathbf G^{-1}\mathbf h}\,q_h^{\frac M2-1}e^{-\frac{Mq_h}{2\bar q_h}}\,
      \exp\!\big[\tfrac{\tilde g}{2}(\mathbf t\cdot\mathbf e(\mathbf h,q_h))^2\big], \nonumber\\
  &\text{(E3)}\quad \mathbf K=\E[\mathbf e\mathbf e^\top],\quad \mathbf H=\E[\mathbf h\mathbf h^\top],\quad\langle q_h\rangle=\E q_h, \nonumber\\
  &\text{(E4)}\quad (\sigma_A^2\mathbf K+\kappa/P)\,\mathbf t=\mathbf y,\qquad \bar{\mathbf f}=\sigma_A^2\mathbf K\mathbf t, \nonumber\\
  &\text{(E5)}\quad p(w_0,q)\propto q^{\frac{d-3}{2}}\exp[-\beta(w_0^2+q)]
      \exp\Big[-\tfrac{r_B}{2}\mathbf c^\top \mathbf G^{-1}(\mathbf G-\mathbf H)\mathbf G^{-1}\mathbf c
      -\tfrac{r_B}{2}\tfrac{M(\bar q_h-\langle q_h\rangle)}{\bar q_h^2}P_\perp\Big], \nonumber\\
  &\text{(E6)}\quad \mathbf J=\E[\mathbf c\mathbf c^\top], \nonumber\\
  &\text{(E7)}\quad \bar q_h=\sigma_B^2\big(\E\|\erf(w\cdot x)\|_x^2-\tr \mathbf J\big), \nonumber\\
  &\text{(E8)}\quad M=\bar q_h^{\,2}\big/\sigma_B^4\big(\E_{w,w'}k(w,w')^2-\|\mathbf J\|_F^2\big). \nonumber
\end{align}

%% ================================================================
\section{The diagonal, Gaussian limit}\label{sec:reduce}
%% ================================================================
Three further approximations reduce (E1)--(E8) to three scalar equations.
Define $\gamma(A)=\frac{2}{\sqrt\pi}(1+2\tr A)^{-1/2}$, and write $\gamma_H=\gamma(H)$ and $\gamma_\Sigma=\gamma(\Sigma)$, where $\Sigma$ is
the covariance of a read-in row.

\paragraph{(a) Linearise $\mathbf e$.}
Set $e_n\approx\gamma_H h_n$, with $q_h\to\tr H$.
Then (E2) is Gaussian, $\mathbf K=\gamma_H^2\mathbf H$, and
\begin{equation}
  \mathbf H=\big(\mathbf G^{-1}-\tilde g\gamma_H^2\,\mathbf t\mathbf t^\top\big)^{-1}
     =\mathbf G+\frac{\tilde g\gamma_H^2\,\mathbf G\mathbf t\,\mathbf t^\top \mathbf G}{1-\tilde g\gamma_H^2\,\mathbf t^\top \mathbf G\mathbf t}.
  \label{eq:Hgauss}
\end{equation}

\paragraph{(b) Keep only the diagonal.}
Write $\lambda^H_n=H_{nn}$, and keep $H_{nn}=[G_{nn}^{-1}-\tilde g\gamma_H^2t_n^2]^{-1}$ with $t_n=y_n/(\kappa/P+\sigma_A^2\gamma_H^2\lambda^H_n)$:
\begin{equation}
  \lambda^H_n=\Big[\frac{1}{G_{nn}}-\tilde g\,\frac{\gamma_H^2\,y_n^2}{\big(\kappa/P+\sigma_A^2\gamma_H^2\lambda^H_n\big)^2}\Big]^{-1},
  \qquad n=1,3 .
  \label{eq:lamH}
\end{equation}
Step (b) is separate from step (a).
Even when $\mathbf G$ is diagonal, \eqref{eq:Hgauss} has
\begin{equation}
  H_{13}=\frac{\tilde g\gamma_H^2\,G_{11}G_{33}\,t_1t_3}{1-\tilde g\gamma_H^2(G_{11}t_1^2+G_{33}t_3^2)}\neq0 .
\end{equation}

\paragraph{(c) Gaussian read-in.}
Take $w\sim\mathcal N(0,\Sigma)$, with eigenvalue $\lambda^\Sigma$ along $v_*$.
Keep only the $c_1^2$ term of (E5), with $D_w\to1+2\tr\Sigma$, so that $c_1^2=\gamma_\Sigma^2w_0^2$.
Then $J_{11}=\gamma_\Sigma^2\lambda^\Sigma$ and, from $\E w_0^6=15(\lambda^\Sigma)^3$,
$J_{33}=\frac{40(\lambda^\Sigma)^3}{\pi(1+2\tr\Sigma)^3}$.
The precision along $v_*$ is $2\beta+r_B\gamma_\Sigma^2\mathcal V_{11}$, with
$\mathcal V_{11}=(G_{11}-H_{11})/G_{11}^2$.
At $\sigma_B=\sigma_w=1$ this is
\begin{equation}
  \lambda^\Sigma=\Big[d+\frac{N_2}{N_1}\,\frac{1}{\gamma_\Sigma^2}\,
  \frac{\gamma_\Sigma^2\lambda^\Sigma-\lambda^H_1}{(\lambda^\Sigma)^2}\Big]^{-1}.
  \label{eq:lamS}
\end{equation}
Equations \eqref{eq:lamH} (for $n=1,3$, with $G_{nn}=J_{nn}$) and \eqref{eq:lamS} form a closed system for
$(\lambda^H_1,\lambda^H_3,\lambda^\Sigma)$.

%% ================================================================
\section{Infinite first layer is not a two-layer network}\label{sec:limit}
%% ================================================================
By \eqref{eq:oneW}, $N_1$ enters only through $r_B=\sigma_B^2N_2/N_1$.
At $N_1=\infty$ the term $\frac{r_B}{2}\langle\Phi\circ W_j|\mathcal V|\Phi\circ W_j\rangle$ in \eqref{eq:oneW} vanishes: (E5) becomes the prior, and in the diagonal limit \eqref{eq:lamS} gives $\lambda^\Sigma=\sigma_w^2/d$.
What remains is (E2)--(E4), the distribution of $h$ with a \emph{fixed} prior.
That prior has covariance $\mathbf G^0=\sigma_B^2\E_{\rm prior}[\mathbf c\mathbf c^\top]$, perpendicular mean $\bar q^{\,0}_h$, and dimension $M^0$.

\paragraph{Two-layer network in the same variables.}
For $f=\sum_ka_k\erf(w_k\cdot x)$ the same construction gives a single-neuron distribution for $w$:
\begin{equation}
  -\log p(w_0,q)=\beta(w_0^2+q)-\tfrac{d-3}{2}\log q-\frac{\tilde g}{2}\big(\mathbf t\cdot\mathbf c(w_0,q)\big)^2 .
\end{equation}
Because $e_n(h_1,0,\dots;q)=c_n(h_1,q)$, this is (E2) with $\mathbf h=(w_0,0,0,\dots)$, $q_h=q$, and
\begin{equation}
  \mathbf G^{\rm 2L}=\mathrm{diag}\big(\tfrac{\sigma_w^2}{d},0,0,\dots\big),\qquad
  \bar q^{\,\rm 2L}=\sigma_w^2\tfrac{d-1}{d},\qquad M^{\rm 2L}=d-1 .
\end{equation}
Here $q=|w_\perp|^2$ is $\sigma_w^2/d$ times a sum of $d-1$ squared standard normals, so \eqref{eq:gamma} with $M=d-1$ is exact.
So the three-layer equations at $N_1=\infty$ coincide with the two-layer ones if and only if
$(\mathbf G^0,\bar q^{\,0}_h,M^0)=(\mathbf G^{\rm 2L},\bar q^{\,\rm 2L},M^{\rm 2L})$.

\paragraph{The infinite first layer.}
From \eqref{eq:cn}, with $w_0\sim\mathcal N(0,\sigma_w^2/d)$ and $D_w\approx D_0=1+2\sigma_w^2$, and with
$\E w_0^2=\sigma_w^2/d$, $\E w_0^4=3\sigma_w^4/d^2$, $\E w_0^6=15\sigma_w^6/d^3$:
\begin{equation}
  J^0_{11}\simeq\frac{4}{\pi}\frac{\sigma_w^2}{dD_0},\qquad
  J^0_{13}\simeq-\frac{24}{\sqrt6\,\pi}\frac{\sigma_w^4}{d^2D_0^2},\qquad
  J^0_{33}\simeq\frac{40}{\pi}\frac{\sigma_w^6}{d^3D_0^3},\qquad
  \rho_{13}=\frac{J^0_{13}}{\sqrt{J^0_{11}J^0_{33}}}=-\frac{3}{\sqrt{15}} .
\end{equation}
$\mathbf G^0$ is therefore not of the two-layer form.
It carries He$_3$ content $\sqrt{G^0_{33}/G^0_{11}}\sim1/d$, with a fixed negative correlation $\rho_{13}=-0.775$ for every $d$.

\begin{table}[!htbp]
\centering
\small
\begin{tabular}{@{}rrrrrrr@{}}
\toprule
$d$ & $M^0$ & $d-1$ & $\bar q^{\,0}_h/G^0_{11}$ & $\sqrt{G^0_{33}/G^0_{11}}$ & $\mu_3$, 3 layers, $N_1=\infty$ & $\mu_3$, 2 layers\\
\midrule
75  & 88.4  & 74  & 81.1  & 0.0135 & 0.0695  & 0.1224\\
150 & 178.2 & 149 & 163.2 & 0.0069 & $-0.1290$ & $-0.0869$\\
300 &       &     &       &        & $-0.0956$ & $-0.0613$\\
600 &       &     &       &        & $-0.0602$ & $-0.0428$\\
1200&       &     &       &        & $-0.0384$ & $-0.0297$\\
\bottomrule
\end{tabular}
\caption{Numerical solutions of (E1)--(E8) at $r_B=0$, with $P=3000$, $N_2=1600$, $\chi=10$, $\kappa=0.1$, $\varepsilon=0.074$,
$\sigma_A=\sigma_B=\sigma_w=1$, $n_{\max}=11$.
``2 layers'' is the two-layer distribution with its prior values $(\sigma_w^2/d,\ \sigma_w^2(d-1)/d,\ d-1)$ replaced by $(G^0_{11},\bar q^{\,0}_h,M^0)$.
It agrees to 7--9 digits with (E1)--(E8) at $r_B=0$ with $G_{13}=G_{33}=0$.
So the remaining difference is the He$_3$ content $G^0_{13},G^0_{33}$ of the infinite first layer.}
\label{tab:limit}
\end{table}

\paragraph{Where the He$_3$ eigenvalue of $H$ comes from.}
\begin{itemize}
  \item \emph{Diagonal, Gaussian limit.} By \eqref{eq:lamH}, the only source of $\lambda^H_3$ is $1/G_{33}$, which is inherited from the first layer.
    The two-layer network in the same limit has $G_{33}=0$, hence $\lambda^H_3=0$.
    So every He$_3$ quantity differs between the models: $\lambda^H_3$ is nonzero for one and zero for the other.
  \item \emph{Exact $\mathbf e$.} By \eqref{eq:Psi}, $e_3$ also contains a self-generated term,
    $e_3\ni c_3(h_1,q_h)\propto-h_1^3/(1+2q_h)^{3/2}$.
    This is the term the two-layer network uses.
    It has size $\sim G_{11}^{3/2}\sim\sqrt{G_{11}}/d$, the same order as the inherited $\sqrt{G_{33}}\sim\sqrt{G_{11}}/d$.
    The inherited part is correlated with $h_1$ with correlation $\rho_{13}<0$, so the two parts of $e_3$ partly cancel.
    The models therefore differ in $\mu_3$ already at the order of $e_3$ itself.
  \item $M^0\neq d-1$ and $\bar q^{\,0}_h/G^0_{11}\neq d-1$ (Table~\ref{tab:limit}) contribute further, through the distribution of $q_h$.
\end{itemize}


%% ================================================================
\section{Approximations}
%% ================================================================
\begin{enumerate}
  \item Dataset-averaged likelihood.
  \item Mean-field decoupling of the quartic terms in \eqref{eq:S}, keeping the $O(\chi^2)$ part of $\mathcal T$.
  \item Restriction to the subspace spanned by $\varphi_n$, $n\in S$; everything outside it enters only through $q_h$ and $P_\perp$.
  \item At fixed $x_0$, $h_\perp$ is Gaussian over $x_\perp$, which gives \eqref{eq:Psi}.
  \item The mean-and-variance approximation \eqref{eq:gamma}--\eqref{eq:M} for the distribution of $q_h$; the covariance \eqref{eq:Gperp} of $h_\perp$; and the equal-eigenvalue treatment of $\mathcal V$ outside the subspace.
\end{enumerate}

%% ================================================================
\appendix
\section{The exact posterior sampler and its fast implementation}\label{sec:sampler}
%% ================================================================
The equations of state are compared with samples from the exact posterior of \eqref{eq:net} on a fixed training set
$\{x_\mu,y_\mu\}_{\mu=1}^P$, i.e.\ \eqref{eq:S0} with $P\|f-y\|_x^2$ replaced by $\sum_\mu(f(x_\mu)-y_\mu)^2$
(code: \texttt{milestones/fcn3\_collapsed\_hmc}).

\paragraph{Moves.}
Each iteration leaves the posterior invariant.
\begin{enumerate}
  \item HMC on $(W,B)$ with the readout $a$ integrated out. Given $(W,B)$, $a$ is Gaussian, so the collapsed potential
    $U(W,B)$ needs one $N_2\times N_2$ Cholesky factorisation of $M=G/v+I/s_a^2$, with $G=F_2^\top F_2$, $F_2=\erf(\erf(XW^\top)B^\top)$, $v=\kappa/\chi$, $s_a^2=\sigma_A^2/(N_2\chi)$.
  \item $a\sim p(a\mid W,B)$, drawn exactly.
  \item Every second iteration, a sweep of per-neuron jump moves.
    A middle neuron $k$ moves $(B_k,a_k)$ jointly, with $a_k$ integrated out.
    A read-in neuron $j$ moves $(w_j,B_{:,j})$ jointly, in the coordinates (aligned component $w_0$, perpendicular norm $q$) of Section~\ref{sec:reduce}.
    Each is a multiple-try independence sampler whose proposal depends only on the other neurons.
  \item The exact sign symmetries $(w_j,B_{:,j})\to-(w_j,B_{:,j})$ and $(B_k,a_k)\to-(B_k,a_k)$.
\end{enumerate}

\paragraph{Fast implementation.}
The fast path (\texttt{fcn3\_fast.py}, the default since 2026-11-09) changes how these moves are computed, not the moves themselves.
\begin{itemize}
  \item \emph{HMC gradient in float32.} A leapfrog step built from any deterministic force $F(q)$,
    \[p\to p+\epsilon F(q),\qquad q\to q+\epsilon p,\]
    is a composition of shears. It is therefore volume preserving and, with the momentum flip, time reversible, whether or not $F=-\nabla U$.
    HMC stays exact if the Metropolis test uses the exact $U$.
    So the leapfrog uses a float32 gradient, and the float64 energy is evaluated only at the two ends of each trajectory.
    On the RTX~5090, float64 runs at $1/64$ of the float32 rate, and the float64 Cholesky factorisation and inverse took about 90\% of the HMC time.
  \item \emph{Sweeps without host synchronisation.}
    One neuron update is captured in a CUDA graph and replayed $N_1$ or $N_2$ times.
    The random numbers of a sweep are drawn in advance from the sampler's generator.
    The choice among candidates is a Gumbel argmax, which has the same law as sampling from the softmax of the candidate weights.
    The selected candidate is written back unconditionally; candidate 0 is the current state.
  \item \emph{Fused kernels for the read-in move.}
    For each candidate $c$, the predictions $\sum_k a_k\erf(h^0_k+\phi_c B_{kc})$ are summed without forming the
    (candidates $\times P\times N_2$) array. Here $h^0=H-\erf(Xw_j)B_{:,j}^\top$ is the hidden pre-activation with neuron $j$ removed.
    The conditional mean of $B_{:,j}$, which enters only the proposal, is computed in bfloat16.
    Any deterministic function of the rest of the state is a valid proposal, so this precision affects the acceptance rate only.
    All candidate energies keep the reference precision: float32 data terms, accumulated in float64.
\end{itemize}

\paragraph{Validation.}
\begin{enumerate}
  \item The fused kernels agree with direct float64 evaluation to $\sim10^{-7}$ (relative).
    The float32 gradient agrees with the float64 one to $8\times10^{-7}$.
    The graphed sweep makes bit-identical decisions to the eager one.
  \item On a tiny network ($d=5$, $P=12$, $N_1=4$, $N_2=3$; $2\times10^4$ iterations each), three samplers of the same posterior agree:
    reference HMC, fast HMC alone, and fast jump moves alone (no HMC).
    Over eight statistics (moments of $w_0$, $|w_\perp|$, $B$, $a$ and of $f$ at test points), the 16 comparisons give $\max|z|=1.81$ (batch means).
  \item Table~\ref{tab:fast} compares full chains at the $\kappa_0=0.015$ ray point.
\end{enumerate}

\begin{table}[h]
\centering
\begin{tabular}{lccc}
\toprule
 & fast, prior start & reference, prior start & reference, aligned start\\
\midrule
$d\langle w_0^2\rangle$ & $7.07\pm0.01$ & $7.14\pm0.02$ & $7.43\pm0.02$\\
$L_1$ & $1.029$ & $1.029$ & $1.030$\\
$L_3$ & $0.5462$ & $0.5455$ & $0.5486$\\
$H_{11}$ & $0.0348$ & $0.0350$ & $0.0359$\\
$H_{33}$ & $0.0036$ & $0.0036$ & $0.0036$\\
read-in mass at $|w_0|>0.5$ & $0.0185$ & $0.0188$ & $0.0200$\\
HMC acceptance & $0.82$ & $0.73$ & $0.65$\\
wall time (1000 iterations) & 452\,s & 2938\,s & 2541\,s\\
\bottomrule
\end{tabular}
\caption{Fast and reference samplers at $d=150$, $P=3000$, $N_1=N_2=1500$, $\chi=1$, $\kappa=0.015$, $\varepsilon=-0.3\sqrt6$,
$\sigma_A^2=8/1500$, $\sigma_B=1.2554$, $\sigma_w^2=1/2$. Each chain ran 1000 iterations with a burn-in of 300.
Errors are batch means over 10 batches; they underestimate the spread between chains, since the two reference chains differ in $d\langle w_0^2\rangle$ by 0.29.
The fast chain shared its GPU with another job; the reference chains did not.}
\label{tab:fast}
\end{table}

\paragraph{Cost.}
At the point of Table~\ref{tab:fast}, on an otherwise idle RTX~5090, one sweep (middle and read-in layers) takes $0.47$\,s and one HMC trajectory $0.07$\,s, against $4.2$\,s and $0.29$\,s for the reference implementation.
With a sweep every second iteration this is $0.31$\,s per iteration, so about 5 minutes per 1000 iterations instead of 49.
The read-in move costs $O(P N_2 n_c)$ per neuron, with $n_c=16$ candidates. A sweep therefore scales as $P N_1 N_2$, and doubling $N$ at fixed $P$ makes it about four times slower.
The remaining cost is set mainly by the $P N_2 n_c$ erf evaluations per read-in neuron.
A Taylor-series evaluation in $\phi_c B_{kc}$ with a rigorous remainder bound was tried and was not faster.
\texttt{run\_fcn3.py --no-fast} runs the reference implementation.

\end{document}
